\documentclass[11pt, letterpaper]{article}

% --- Idioma y codificación ---
\usepackage[spanish]{babel}
\usepackage[utf8]{inputenc}
\usepackage[T1]{fontenc}
\usepackage{lmodern}

% --- Matemáticas ---
\usepackage{amsmath, amssymb, amsthm}

% --- Formato ---
\usepackage[margin=2.5cm]{geometry}
\usepackage{parskip}
\usepackage{booktabs}
\usepackage{hyperref}
\hypersetup{colorlinks=true, linkcolor=blue!70!black, urlcolor=blue!70!black}

% --- Entornos de teoremas ---
\theoremstyle{definition}
\newtheorem{definition}{Definición}[section]
\theoremstyle{plain}
\newtheorem{theorem}{Teorema}[section]
\newtheorem{lemma}[theorem]{Lema}
\newtheorem{corollary}[theorem]{Corolario}
\theoremstyle{remark}
\newtheorem{remark}{Observación}[section]

\title{%
  T05 --- Demostración formal:\\[4pt]
  Inserción con rehash es $O(1)$ amortizado
}
\author{Curso Linux--C--Rust --- B15 Estructuras de datos en C}
\date{}

\begin{document}
\maketitle
\tableofcontents

% ===================================================================
\section{Modelo}
% ===================================================================

\begin{definition}
Tabla hash dinámica con capacidad inicial~$m_0$, factor de
crecimiento $\times 2$, umbral de carga $\alpha_{\max}$ (constante).
Cuando una inserción haría que $n > \alpha_{\max} \cdot m$, se crea
una nueva tabla de capacidad~$2m$ y se reinsertan los $n$ elementos
(rehash completo).
\end{definition}

\textbf{Costos unitarios:}
\begin{itemize}
  \item Inserción sin resize: $c_{\text{ins}} = 1$.
  \item Reinserción durante rehash: $c_{\text{re}} = 1$.
  \item Resize con $n$ elementos: $c_{\text{resize}} = n$.
\end{itemize}

\textbf{Objetivo.} Demostrar que $n$ inserciones consecutivas
(empezando con tabla vacía) cuestan $O(n)$ en total.

% ===================================================================
\section{Demostración 1: método agregado}
% ===================================================================

\begin{theorem}\label{thm:aggregate}
El costo total de $n$ inserciones es a lo sumo $3n$.
\end{theorem}

\begin{proof}
Partimos de una tabla vacía con capacidad~$m_0$. Sea
$n_k = \alpha_{\max} \cdot 2^k m_0$ el número de elementos que
dispara el $k$-ésimo resize.

\textbf{1) Inserciones ordinarias.} Cada una de las $n$ inserciones
cuesta~1:
\[
  C_{\text{ins}} = n
\]

\textbf{2) Rehashes.} Sea $K$ el número total de resizes. El
$k$-ésimo resize reinserta $n_k$ elementos:
\[
  C_{\text{rehash}}
  = \sum_{k=0}^{K} n_k
  = \sum_{k=0}^{K} \alpha_{\max} \cdot 2^k m_0
  = \alpha_{\max} m_0 \cdot \frac{2^{K+1} - 1}{2 - 1}
  = \alpha_{\max} m_0 (2^{K+1} - 1)
\]

Como $n_K = \alpha_{\max} \cdot 2^K m_0 \leq n$, tenemos
$2^K \leq n/(\alpha_{\max} m_0)$:
\[
  C_{\text{rehash}}
  \leq \alpha_{\max} m_0
    \left(\frac{2n}{\alpha_{\max} m_0} - 1\right)
  = 2n - \alpha_{\max} m_0 < 2n
\]

\textbf{Total:}
\[
  C_{\text{total}} = C_{\text{ins}} + C_{\text{rehash}}
  < n + 2n = 3n
\]

Costo amortizado por inserción:
$\hat{c} = C_{\text{total}}/n < 3 = O(1)$.
\end{proof}

% ===================================================================
\section{Demostración 2: método del banquero (créditos)}
% ===================================================================

\begin{theorem}\label{thm:banker}
Si cada inserción paga 3 unidades de costo, el saldo de créditos
nunca es negativo.
\end{theorem}

\begin{proof}
Cada inserción usa sus 3 unidades así:
\begin{itemize}
  \item \textbf{1 unidad}: paga la inserción misma.
  \item \textbf{2 unidades}: se depositan como crédito, asociadas al
    elemento insertado.
\end{itemize}

\textbf{Invariante de crédito.} Después de cada operación, cada
elemento insertado desde el último resize tiene al menos 2 créditos.

\textbf{Verificación al momento del resize.} Cuando se dispara un
resize, la tabla tiene $n_{\text{cur}} = \alpha_{\max} \cdot m$
elementos. De estos, $n_{\text{cur}}/2$ fueron insertados desde el
resize anterior (cuando la tabla tenía $\alpha_{\max} \cdot m/2$
elementos y se expandió a capacidad~$m$, quedando con factor
$\alpha_{\max}/2$).

Créditos disponibles:
\[
  \text{créditos} = 2 \cdot \frac{n_{\text{cur}}}{2} = n_{\text{cur}}
\]

Costo del resize:
\[
  \text{costo resize} = n_{\text{cur}}
\]

Los créditos \textbf{cubren exactamente} el costo. Cada nuevo
elemento paga por sí mismo y por un viejo.

Después del resize, todos los créditos se consumieron, pero el
invariante se restablece conforme llegan nuevas inserciones con
2~créditos cada una. El saldo nunca es negativo. \checkmark
\end{proof}

% ===================================================================
\section{Demostración 3: método del potencial}
% ===================================================================

\begin{theorem}\label{thm:potential}
Con la función de potencial $\Phi = 2n - \alpha_{\max} \cdot m$, el
costo amortizado de cada inserción es exactamente~3.
\end{theorem}

\begin{proof}
Definimos $\Phi = 2n - \alpha_{\max} \cdot m$ donde $n$ es el número
de elementos y $m$ la capacidad actual.

\textbf{Propiedades del potencial:}
\begin{itemize}
  \item Tras un resize: $n = \alpha_{\max} \cdot m/2$, así que
    $\Phi = 2 \cdot \alpha_{\max} m/2 - \alpha_{\max} m = 0$.
  \item Justo antes del siguiente resize: $n = \alpha_{\max} \cdot m$,
    así que $\Phi = 2\alpha_{\max} m - \alpha_{\max} m = n$.
  \item $\Phi \geq 0$ siempre: entre resizes, $n$ crece desde
    $\alpha_{\max} m/2$ hasta $\alpha_{\max} m$, y
    $\Phi = 2n - \alpha_{\max} m$ va de~0 a~$n$.
\end{itemize}

\textbf{Caso 1: inserción sin resize.} $n$ aumenta en~1, $m$ no
cambia:
\[
  \hat{c} = c_{\text{real}} + \Delta\Phi
  = 1 + [2(n{+}1) - \alpha_{\max} m] - [2n - \alpha_{\max} m]
  = 1 + 2 = 3
\]

\textbf{Caso 2: inserción con resize.} Justo antes:
$n = \alpha_{\max} m$, $\Phi_{\text{antes}} = n$. El resize cambia
$m \to 2m$, luego se inserta: $n \to n + 1$.
\[
  \Phi_{\text{después}}
  = 2(n{+}1) - \alpha_{\max} \cdot 2m
  = 2(n{+}1) - 2n = 2
\]

Costo real: $c_{\text{real}} = n + 1$ (reinsertar $n$ + insertar el
nuevo):
\[
  \hat{c} = c_{\text{real}} + \Delta\Phi
  = (n + 1) + (2 - n) = 3
\]

\textbf{En ambos casos, $\hat{c} = 3$.}

Costo total real:
\[
  \sum_{i=1}^{n} c_i
  = \sum_{i=1}^{n} \hat{c}_i - \Phi_n + \Phi_0
  \leq 3n - \Phi_n + 0 \leq 3n \qedhere
\]
\end{proof}

% ===================================================================
\section{Demostración 4: crecimiento aditivo es
  \texorpdfstring{$\Theta(n)$}{Theta(n)} amortizado}
% ===================================================================

\begin{theorem}\label{thm:additive}
Si la tabla crece sumando una constante~$c$ (en vez de multiplicar),
el costo amortizado por inserción es $\Theta(n)$, no $O(1)$.
\end{theorem}

\begin{proof}
Con crecimiento aditivo ($m \to m + c$), los resizes ocurren cuando
$n$ alcanza $\alpha_{\max}(m_0 + kc)$ para $k = 0, 1, 2, \ldots$

El $k$-ésimo resize reinserta $\alpha_{\max}(m_0 + kc)$ elementos.
El número de resizes hasta insertar $n$ elementos es
$K = \Theta(n/c)$.

Costo total de los rehashes:
\[
  C_{\text{rehash}}
  = \sum_{k=0}^{K} \alpha_{\max}(m_0 + kc)
  = \alpha_{\max}\!\left(Km_0 + c \cdot \frac{K(K{+}1)}{2}\right)
  = \Theta(K^2)
  = \Theta(n^2/c)
\]

Costo amortizado:
$\hat{c} = \Theta(n^2/c) / n = \Theta(n/c) = \Theta(n)$.
\end{proof}

\begin{remark}
La diferencia es suma aritmética $\sum k = \Theta(K^2)$ vs suma
geométrica $\sum 2^k = \Theta(2^K)$. Por esto las tablas hash usan
crecimiento multiplicativo.
\end{remark}

% ===================================================================
\section{Demostración 5: shrink con histéresis preserva $O(1)$}
% ===================================================================

\begin{theorem}\label{thm:shrink}
Si la tabla encoge a la mitad cuando
$\alpha < \alpha_{\min}$ con
$\alpha_{\min} \leq \alpha_{\max}/4$, entonces una secuencia
intercalada de $n$ inserciones y eliminaciones tiene costo amortizado
$O(1)$.
\end{theorem}

\begin{proof}
Extendemos el potencial:
\[
  \Phi = \begin{cases}
    2n - \alpha_{\max} \cdot m
      & \text{si } n \geq \alpha_{\max} \cdot m / 2 \\
    \alpha_{\min} \cdot m - n
      & \text{si } n < \alpha_{\max} \cdot m / 2
  \end{cases}
\]

\textbf{Propiedad clave:} $\Phi = 0$ cuando
$n = \alpha_{\max} m / 2$ (punto medio), y $\Phi \geq 0$ siempre.
\begin{itemize}
  \item Antes de un grow: $n = \alpha_{\max} m$, $\Phi = n$. El
    costo del grow ($= n$) se paga con el potencial.
  \item Antes de un shrink: con
    $\alpha_{\min} \leq \alpha_{\max}/4$, se acumulan suficientes
    eliminaciones entre el punto medio y el umbral para financiar el
    rehash.
\end{itemize}

\textbf{Anti-thrashing.} Después de un grow,
$\alpha = \alpha_{\max}/2$. Para disparar un shrink, $\alpha$ debe
bajar hasta $\alpha_{\min} \leq \alpha_{\max}/4$, requiriendo al
menos $\alpha_{\max} m / 4$ eliminaciones. Cada una aporta potencial.
Después de un shrink, $\alpha = 2\alpha_{\min} \leq \alpha_{\max}/2$,
lejos de ambos umbrales. No hay oscilación.
\end{proof}

% ===================================================================
\section{Verificación numérica}
% ===================================================================

Tabla con $m_0 = 8$, $\alpha_{\max} = 0.75$, factor $\times 2$:

\begin{table}[h]
\centering
\begin{tabular}{cccccc}
\toprule
Inserción $i$ & $n$ & ¿Resize? & Costo real & Créd.\ usados
  & Créd.\ rest. \\
\midrule
1--6   & 1--6  & No  & 1 c/u   & 0  & $6 \times 2 = 12$ \\
7      & 7     & Sí ($8 \to 16$)  & $6{+}1=7$   & 6  & $12{-}6{+}2=8$ \\
8--12  & 8--12 & No  & 1 c/u   & 0  & $8{+}5{\times}2=18$ \\
13     & 13    & Sí ($16 \to 32$) & $12{+}1=13$ & 12 & $18{-}12{+}2=8$ \\
14--24 & 14--24 & No & 1 c/u   & 0  & $8{+}11{\times}2=30$ \\
25     & 25    & Sí ($32 \to 64$) & $24{+}1=25$ & 24 & $30{-}24{+}2=8$ \\
\bottomrule
\end{tabular}
\end{table}

Costo total para 25 inserciones: $25 + 6 + 12 + 24 = 67$.
Cota del teorema: $3 \times 25 = 75$. \checkmark \quad ($67 < 75$).

Costo amortizado observado: $67/25 = 2.68 < 3$. \checkmark

% ===================================================================
\section{Resumen de resultados}
% ===================================================================

\begin{table}[h]
\centering
\begin{tabular}{llp{5.5cm}}
\toprule
\textbf{Resultado} & \textbf{Técnica} & \textbf{Clave} \\
\midrule
$C_{\text{total}} < 3n$
  & Método agregado
  & Serie geométrica $\sum 2^k < 2n$ \\
3 unidades/ins.\ bastan
  & Método del banquero
  & 2 créditos/elem.; cada nuevo paga por sí mismo y un viejo \\
$\hat{c} = 3$ exacto
  & Potencial $\Phi = 2n {-} \alpha_{\max}m$
  & $\Phi = 0$ tras resize, $\Phi = n$ antes \\
Crecimiento aditivo: $\Theta(n)$
  & Suma aritmética vs geométrica
  & $\sum k = \Theta(K^2)$ vs $\sum 2^k = \Theta(2^K)$ \\
Shrink con histéresis: $O(1)$
  & Potencial extendido
  & $\alpha_{\min} \leq \alpha_{\max}/4$ evita thrashing \\
\bottomrule
\end{tabular}
\end{table}

\end{document}
